\documentclass[a4paper,UKenglish,cleveref, autoref, thm-restate]{lipics-v2021}
\bibliographystyle{plainurl}
\title{Anders: Modal Homotopy Type System for Computable Modal Homotopy Type Theory and Differential Geometry}
\titlerunning{Anders 5.5.0}
\author{Maksym Sokhatskyi}{Groupoid Infinity}{maxim@synrc.com}{https://orcid.org/0000-0001-7127-8796}{}
\authorrunning{Maksym Sokhatskyi}
\Copyright{Maksym Sokhatskyi}

\begin{CCSXML}
<ccs2012>
<concept>
<concept_id>10003752.10003753.10003754.10003733</concept_id>
<concept_desc>Theory of computation ~ Lambda calculus</concept_desc>
<concept_significance>300</concept_significance>
</concept>
</ccs2012>
\end{CCSXML}

\ccsdesc[300]{Theory of computation~Lambda calculus}
\keywords{Homotopy Type Theory, Differential Geometry, Cohesive Type Theory, Simplicial Type Theory, Modal Type Theory, Foundations of Mathematics}

\nolinenumbers
\hideLIPIcs

\frenchspacing

\newcommand{\tabstyle}[0]{\scriptsize\ttfamily\fontseries{l}\selectfont}

\newcommand{\anders}[0]{\textbf{Anders}}

\usepackage[utf8]{inputenc}
\DeclareUnicodeCharacter{2071}{\textsuperscript{i}}
\DeclareUnicodeCharacter{2080}{\ensuremath{_0}}
\DeclareUnicodeCharacter{2081}{\ensuremath{_1}}
\DeclareUnicodeCharacter{2082}{\ensuremath{_2}}
\DeclareUnicodeCharacter{207B}{\ensuremath{^{-}}}
\DeclareUnicodeCharacter{02B2}{\textsuperscript{j}}
\DeclareUnicodeCharacter{03C6}{\ensuremath{\varphi}}
\DeclareUnicodeCharacter{21A6}{\ensuremath{\mapsto}}
\DeclareUnicodeCharacter{3008}{\ensuremath{\langle}}
\DeclareUnicodeCharacter{3009}{\ensuremath{\rangle}}
\DeclareUnicodeCharacter{266D}{\ensuremath{\flat}}
\DeclareUnicodeCharacter{2111}{\ensuremath{\Im}}
\DeclareUnicodeCharacter{03B9}{\ensuremath{\iota}}
\DeclareUnicodeCharacter{2227}{\ensuremath{\wedge}}
\DeclareUnicodeCharacter{2228}{\ensuremath{\vee}}
\DeclareUnicodeCharacter{03A0}{\ensuremath{\Pi}}
\DeclareUnicodeCharacter{03A3}{\ensuremath{\Sigma}}
\DeclareUnicodeCharacter{2032}{\ensuremath{\prime}}
\DeclareUnicodeCharacter{2212}{\ensuremath{-}}
\DeclareUnicodeCharacter{03BB}{\ensuremath{\lambda}}
\DeclareUnicodeCharacter{03C8}{\ensuremath{\psi}}
\DeclareUnicodeCharacter{2605}{\ensuremath{\star}}
\DeclareUnicodeCharacter{266F}{\ensuremath{\sharp}}
\DeclareUnicodeCharacter{03BC}{\ensuremath{\mu}}

\begin{document}

\maketitle

\begin{abstract}
Derivative work of \textbf{cubicaltt} with \textbf{cohesivett} and \textbf{stt} features.
\end{abstract}

\tableofcontents

\newpage
\section{Introduction}
\label{sec:typesetting-summary}

\textbf{Anders} is a Modal HoTT proof assistant based on: classical MLTT-80 \cite{MLTT80}
with 0, 1, 2, W types; CCHM \cite{CCHM} in CHM \cite{CHM} flavour as cubical type system with
hcomp/transp operations; HTS \cite{HTS} strict equality on pretypes;
infinitisemal \cite{deRham} modality primitives for differential geometry purposes; Disc and Coequalizer primitives.
We tend not to touch general recursive higher inductive schemes,
instead we will try to express as much HIT as possible through Hub Spoke Disc, Coequalizer, and W
primitives built into type checker core in the style of HoTT/Coq homotopy library and Three-HIT theorem \cite{ThreeHIT}.

The HTS language proposed by Voevodsky exposes two different presheaf models of type theory:
the inner one is homotopy type system presheaf that models HoTT and the outer one is
traditional Martin-Löf type system presheaf that models set theory with UIP.
The motivation behind this doubling is to have an ability to express semisemplicial types.
Theoretical work on merging inner and outer languages was continued in 2LTT \cite{2LTT}.

$\mathbf{Installation}$. While we are on our road to Lean-like tactic language, currently we are at the stage of
regular cubical HTS type checker with CHM-style primitives. You may try it from Github
sources: groupoid/anders\footnote{\url{https://github.com/groupoid/anders/}} or install
through OPAM package manager. Main commands are \textbf{check} (to check a program)
and \textbf{repl} (to enter the proof shell).

\begin{table}[ht!]
\centering
\begin{tabular}{l}
\$ opam install anders
\end{tabular}
\end{table}

Anders is fast, idiomatic and educational. We carefully draw the favourite Lean-compatible
syntax to fit 300 LOC in Menhir. The CHM kernel is 2K LOC. Whole Anders compiles under 3
seconds and checks all the base library under 1 minute [i5-12400]. Anders proof assistant
as Homotopy Type System comes with its own Homotopy Library\footnote{\url{https://anders.groupoid.space/library/}}.

\section{Syntax}

The syntax resembles original syntax of the reference CCHM type checker cubicaltt,
is slightly compatible with Lean syntax and contains the full set of Cubical Agda \cite{CubicalAgda}
primitives (except generic higher inductive schemes).

Here is given the mathematical pseudo-code notation of the language
expressions that come immediately after parsing. The core syntax definition of HTS language
corresponds to the type defined in OCaml module:

\begin{table}[ht!]
\centering
\begin{tabular}{rl}
 cosmos :=&$\textbf{U}_j \ |\ \textbf{V}_k$ \\
    var :=&$\textbf{var}\ name\ |\ \textbf{hole}$ \\
     pi :=&$\Pi\ name\ E\ E\ |\ \lambda\ name\ E\ E\ |\ E\ E$ \\
  sigma :=&$\Sigma\ name\ E\ E\ |\ (E,E)\ |\ E.1\ |\ E.2$ \\
      0 :=&$\textbf{0}\ |\ \textbf{ind$_0$}\ \textrm{E}\ \textrm{E}\ \textrm{E}$ \\
      1 :=&$\textbf{1}\ |\ \star\ |\ \textbf{ind$_1$}\ \textrm{E}\ \textrm{E}\ \textrm{E}$ \\
      2 :=&$\textbf{2}\ |\ 0_2\ |\ 1_2\ |\ \textbf{ind$_2$}\ \textrm{E}\ \textrm{E}\ \textrm{E}$ \\
      W :=&$\textbf{W}\ \textrm{ident}\ \textrm{E}\ \textrm{E}\ |\ \textbf{sup}\ \textrm{E}\ \textrm{E}\ |\ \textbf{ind$_W$}\ \textrm{E}\ \textrm{E}$ \\
     id :=&$\textbf{Id}\ E\ |\ \textbf{ref}\ E\ |\ \textbf{id$_J$}\ E$ \\
   path :=&$\textbf{Path}\ E\ |\ E^i\ |\ E\ @\ E$ \\
      I :=&$\textbf{I}\ |\ 0\ |\ 1\ |\ E\ \bigvee\ E\ |\ E\ \bigwedge\ E\ |\ \neg E$ \\
   part :=&$\textbf{Partial}\ E\ E\ |\ \textbf{[}\ (E=I) \rightarrow E, ...\ \textbf{]}$ \\
    sub :=&$\textbf{inc}\ E\ |\ \textbf{ouc}\ E\ |\ E\ \textbf{[}\ I\ \mapsto\ E\ \textbf{]}$ \\
    kan :=&$\textbf{transp}\ E\ E\ |\ \textbf{hcomp}\ E$ \\
   glue :=&$\textbf{Glue}\ E\ |\ \textbf{glue}\ E\ |\ \textbf{unglue}\ E\ E$ \\
     Im :=&$\textbf{Im}\ \mathrm{E}\ |\ \textbf{Inf}\ \mathrm{E}\ |\ \textbf{Join}\ \mathrm{E}\ |\ \textbf{ind$_{Im}$}\ \mathrm{E}\ \mathrm{E}$ \\
   Flat :=&$\flat\ \mathrm{E}\ |\ \textbf{unit}_{\flat}\ \mathrm{E}\ |\ \textbf{counit}_{\flat}\ \mathrm{E}\ |\ \textbf{ind}_{\flat}\ \mathrm{E}\ \mathrm{E}$ \\
  coequ :=&$\textbf{coequ}\ \mathrm{E}\ |\ \iota_2\ |\ \textbf{resp}\ |\ \textbf{ind}_{coequ}$ \\
   disc :=&$\textbf{disc}\ \mathrm{E}\ |\ \textbf{base}\ |\ \textbf{hub}\ |\ \textbf{spoke}\ |\ \textbf{ind}_{disc}$ \\
\end{tabular}
\end{table}

Further Menhir BNF notation will be used to describe the top-level language E parser.

\begin{table}[ht!]
\centering
\begin{tabular}{rl}
      E :=&$\textrm{cosmos}\ |\ \textrm{var}\ |\ \textrm{MLTT}\ |\ \textrm{CCHM}\ |\ \textrm{Im}\ |\ \textrm{Flat}\ |\ \textrm{HIT}$ \\
   CCHM :=&$\textrm{path}\ |\ \textrm{I}\ |\ \textrm{part}\ |\ \textrm{sub}\ |\ \textrm{kan}\ |\ \textrm{glue}$ \\
   MLTT :=&$\textrm{pi}\ |\ \textrm{sigma}\ |\ \textrm{id}$ \\
    HIT :=&$\textrm{coequ}\ |\ \textrm{disc}$ \\
\end{tabular}
\end{table}

\textbf{Keywords}. The words of a top-level language, file or repl, consist of keywords or identifiers.
The keywords are following: \texttt{module}, \texttt{where}, \texttt{import}, \texttt{option}, \texttt{def}, \texttt{axiom},
\texttt{postulate}, \texttt{theorem}, \texttt{(}, \texttt{)}, \texttt{[}, \texttt{]}, \texttt{<}, \texttt{>},
\texttt{/}, \texttt{.1}, \texttt{.2}, \texttt{$\Pi$}, \texttt{$\Sigma$}, \texttt{,}, \texttt{$\lambda$},
\texttt{V}, \texttt{$\bigvee$}, \texttt{$\bigwedge$}, \texttt{-}, \texttt{+}, \texttt{@}, \texttt{PathP},
\texttt{transp}, \texttt{hcomp}, \texttt{zero}, \texttt{one}, \texttt{Partial}, \texttt{inc},
\texttt{$\times$}, \texttt{$\rightarrow$}, \texttt{:}, \texttt{:=}, \texttt{$\mapsto$}, \texttt{U},
\texttt{ouc}, \texttt{interval}, \texttt{inductive}, \texttt{Glue}, \texttt{glue}, \texttt{unglue},
\texttt{$\flat$}, \texttt{$\flat$-unit}, \texttt{$\flat$-counit}, \texttt{ind-$\flat$}, \texttt{coequ}, \texttt{$\iota_2$}, \texttt{resp}, \texttt{coequ-ind}, \texttt{disc}, \texttt{base}, \texttt{hub}, \texttt{spoke}, \texttt{disc-ind}.

\textbf{Indentifiers}. Identifiers support UTF-8. Indentifiers couldn't
start with $\texttt{:}$, $\texttt{-}$, $\rightarrow$. Sample identifiers:
$\neg\texttt{-of-}\vee$, $\texttt{1}$$\rightarrow$$\texttt{1}$, $\texttt{is-?}$,
$\texttt{=}$, $\texttt{\$$\sim$]!}$$\texttt{005x}$, $\infty$, $\texttt{x}$$\rightarrow$$\texttt{Nat}$.

\textbf{Modules}. Modules represent files with declarations. More accurate, BNF notation of module consists of imports, options and declarations.
\begin{table}[ht!]
\tabstyle
\begin{tabular}{l}
\textbf{menhir} \\
\ \ \ start <Module.file> file \\
\ \ \ start <Module.command> repl \\
\ \ \ repl : COLON IDENT exp1 EOF | COLON IDENT EOF | exp0 EOF | EOF \\
\ \ \ file : MODULE IDENT WHERE line* EOF \\
\ \ \ path : IDENT \\
\ \ \ line : IMPORT path+ | OPTION IDENT IDENT | declarations \\
\end{tabular}
\end{table}

\textbf{Imports}. The import construction supports file folder structure (without file extensions)
by using reserved symbol / for hierarchy walking.

\textbf{Options}. Each option holds bool value. Language supports following options:
1) girard (enables U : U);
2) pre-eval (normalization cache);
3) impredicative (infinite hierarchy with impredicativity rule);
In Anders you can enable or disable language core types, adjust syntaxes or
tune inner variables of the type checker.

\textbf{Declarations}. Language supports following top level declarations:
1) axiom (non-computable declaration that breakes normalization);
2) postulate (alternative or inverted axiom that can preserve consistency);
3) definition (almost any explicit term or type in type theory);
4) lemma (helper in big game);
5) theorem (something valuable or complex enough).

\begin{table}[ht!]
\tabstyle
\begin{tabular}{l}
\textbf{axiom} isProp (A : U) : U \\
\textbf{def} isSet (A : U) : U := \textbf{$\Pi$} (a b : A) (x y : Path A a b), Path (Path A a b) x y
\end{tabular}
\end{table}

Sample declarations. For example, signature isProp (A : U) of type U could be
defined as normalization-blocking axiom without proof-term or by providing proof-term as definition.

In this example (A : U), (a b : A) and (x y : Path A a b) are called telescopes.
Each telescope consists of a series of lenses or empty. Each lense provides a
set of variables of the same type. Telescope defines parameters of a declaration.
Types in a telescope, type of a declaration and a proof-terms are a language expressions exp1.

\begin{table}[ht!]
\tabstyle
\begin{tabular}{l}
\textbf{menhir} \\
\ \ \ ident : IRREF | IDENT \\
\ \ \ lense : LPARENS ident+ COLON exp1 RPARENS \\
\ \ \ telescope : lense telescope \\
\ \ \ params : telescope | [] \\
\ \ \ declarations : \\
\ \ \ \ \ \ | DEF IDENT params DEFEQ exp1 \\
\ \ \ \ \ \ | DEF IDENT params COLON exp1 DEFEQ exp1 \\
\ \ \ \ \ \ | AXIOM IDENT params COLON exp1
\end{tabular}
\end{table}

\textbf{Expressions}. All atomic language expressions are grouped by four categories:
exp0 (pair constructions), exp1 (non neutral constructions), exp2 (path and pi applcations),
exp3 (neutral constructions).

\begin{table}[ht!]
\tabstyle
\begin{tabular}{ll}
\textbf{menhir} \\
\multicolumn{2}{l}{\ \ \ face : LPARENS IDENT IDENT IDENT RPARENS } \\
\multicolumn{2}{l}{\ \ \ part : face+ ARROW exp1 } \\
\multicolumn{2}{l}{\ \ \ exp0 : exp1 COMMA exp0 | exp1 } \\
\multicolumn{2}{l}{\ \ \ exp1 : LSQ separated(COMMA, part) RSQ } \\
\ \ \ \ \ \ | LAM telescope COMMA exp1   & | PI telescope COMMA exp1 \\
\ \ \ \ \ \ | SIGMA telescope COMMA exp1 & | LSQ IRREF ARROW exp1 RSQ \\
\ \ \ \ \ \ | LT ident+ GT exp1          & | exp2 ARROW exp1 \\
\ \ \ \ \ \ | exp2 PROD exp1             & | exp2 \\
\end{tabular}
\end{table}

The LR parsers demand to define exp1 as expressions that cannot be used (without a parens enclosure)
as a right part of left-associative application for both Path and Pi lambdas.
Universe indicies $U_j$ (inner fibrant), $V_k$ (outer pretypes) and S (outer strict omega)
are using unicode subscript letters that are already processed in lexer.

\begin{table}[ht!]
\tabstyle
\begin{tabular}{llll}
\textbf{menhir} \\
\multicolumn{3}{l}{\ \ \ exp2 : exp2 exp3 | exp2 APPFORMULA exp3 | exp3 } \\
\multicolumn{4}{l}{\ \ \ exp3 : LPARENS exp0 RPARENS LSQ exp0 MAP exp0 RSQ } \\
\ \ \ \ \ \   | HOLE              & | PRE          & | KAN          & | IDJ exp3 \\
\ \ \ \ \ \   | exp3 FST          & | exp3 SND     & | NEGATE exp3  & | INC exp3 \\
\ \ \ \ \ \   | exp3 AND exp3     & | exp3 OR exp3 & | ID exp3      & | REF exp3\\
\ \ \ \ \ \   | OUC exp3          & | PATHP exp3   & | PARTIAL exp3 & | IDENT \\
\multicolumn{3}{l}{\ \ \ \ \ \    | IDENT LSQ exp0 MAP exp0 RSQ }   & | HCOMP exp3 \\
\multicolumn{3}{l}{\ \ \ \ \ \    | LPARENS exp0 RPARENS }          & | TRANSP exp3 exp3 \\
\end{tabular}
\end{table}


\newpage
\section{Semantics}
The idea is to have a unified layered type checker, so you can disbale/enable any MLTT-style inference,
assign types to universes and enable/disable hierachies. This will be done by providing linking API for
pluggable presheaf modules. We selected 5 levels of type checker awareness from universes and pure type
systems up to synthetic language of homotopy type theory. Each layer corresponds to its presheaves with
separate configuration for universe hierarchies.
\begin{table}[ht!]
\tabstyle
\begin{tabular}{rl}
     \textbf{def} & lang :\textbf{U} := \textbf{inductive} \\
               \{ & UNI: cosmos $\rightarrow$ lang \\
               |  & PI: pure lang $\rightarrow$ lang \\
               |  & SIGMA: total lang $\rightarrow$ lang \\
               |  & ID: strict lang $\rightarrow lang$ \\
               |  & PATH: homotopy lang $\rightarrow$ lang \\
               |  & GLUE: glue lang $\rightarrow$ lang \\
               |  & INDUCTIVE: w012 lang $\rightarrow$ lang \\
               \} & \\
\end{tabular}
\end{table}
We want to mention here with homage to its authors all categorical models of dependent type theory:
Comprehension Categories (Grothendieck, Jacobs), LCCC (Seely), D-Categories and CwA (Cartmell),
CwF (Dybjer), C-Systems (Voevodsky), Natural Models (Awodey). While we can build some transports
between them, we leave this excercise for our mathematical components library.
We will use here the Coquand's notation for Presheaf Type Theories in terms of restriction maps.

\subsection{Universe Hierarchies}

Language supports Agda-style hierarchy of universes: prop, fibrant (U), interval pretypes (V) and
strict omega with explicit level manipulation. All universes are bounded with preorder
\begin{equation}
Fibrant_j \prec Pretypes_k
\end{equation}
in which $j,k$ are bounded with equation:
\begin{equation}
j < k.
\end{equation}
Large elimination to upper universes is prohibited. This is extendable to Agda model:

\begin{table}[ht!]
\tabstyle
\begin{tabular}{rl}
     \textbf{def} & cosmos : \textbf{U} := \textbf{inductive} \\
               \{ & fibrant: nat \\
               |  & pretypes: nat \\
               \} & \\
\end{tabular}
\end{table}

The \textbf{anders} model contains only fibrant $U_j$ and pretypes $V_k$ universe hierarchies.

\newpage

\subsection{Dependent Types}

\begin{definition}[Type]
A type is interpreted as a presheaf $A$, a family of sets $A_I$ with restriction maps
$u \mapsto u\ f, A_I \rightarrow A_J$ for $f: J\rightarrow I$. A dependent type
B on A is interpreted by a presheaf on category of elements of $A$: the objects
are pairs $(I,u)$ with $u : A_I$ and morphisms $f: (J,v) \rightarrow (I,u)$ are
maps $f : J \rightarrow$ such that $v = u\ f$. A dependent type B is thus given
by a family of sets $B(I,u)$ and restriction maps $B(I,u) \rightarrow B(J,u\ f)$.
\end{definition}


We think of $A$ as a type and $B$ as a family of presheves $B(x)$ varying $x:A$.
The operation $\Pi(x:A)B(x)$ generalizes the semantics of
implication in a Kripke model.

\begin{definition}[Pi]
An element $w:[\Pi(x:A)B(x)](I)$ is a family of functions $w_f : \Pi(u:A(J))B(J,u)$
for $f : J \rightarrow I$ such that $(w_f u)g=w_{f\ g}(u\ g)$ when $u:A(J)$ and $g:K\rightarrow J$.
\end{definition}
\begin{table}[ht!]
\tabstyle
\begin{tabular}{rl}
  \textbf{def} & pure (lang : U) : \textbf{U} := \textbf{inductive} \\
            \{ & pi: name $\rightarrow$ nat $\rightarrow$ lang $\rightarrow$ lang $\rightarrow$ pure lang \\
             | & lambda: name $\rightarrow$ nat $\rightarrow$ lang $\rightarrow$ lang \\
             | & app: lang $\rightarrow$ lang \\
            \} & \\
\end{tabular}
\end{table}

\begin{definition}[Sigma]
The set $\Sigma(x:A)B(x)$ is the set of pairs $(u,v)$ when $u:A(I),v:B(I,u)$ and restriction map $(u,v)\ f=(u\ f,v\ f)$.
\end{definition}
\begin{table}[ht!]
\tabstyle
\begin{tabular}{rl}
  \textbf{def} & total (lang : U) : \textbf{U} := \textbf{inductive} \\
            \{ & sigma: name $\rightarrow$ lang $\rightarrow$ total lang \\
             | & pair: lang $\rightarrow$ lang \\
             | & fst: lang \\
             | & snd: lang \\
            \} & \\
\end{tabular}
\end{table}

The presheaf with only Pi and Sigma is called \textbf{MLTT-72} \cite{MLTT72}. Its internalization in \textbf{anders} is as follows:

\begin{table}[ht!]
\tabstyle
\begin{tabular}{rl}
\textbf{def} & MLTT-72 (A : U) (B : A $\rightarrow$ U) : U := $\Sigma$ \\
             & ($\Pi$-form$_1$ : U) \\
             & ($\Pi$-ctor$_1$ : Pi A B $\rightarrow$ Pi A B) \\
             & ($\Pi$-elim$_1$ : Pi A B $\rightarrow$ Pi A B) \\
             & ($\Pi$-comp$_1$ : (a : A) (f : Pi A B), $\Pi$-elim$_1$ ($\Pi$-ctor$_1$ f) a = f a) \\
             & ($\Pi$-comp$_2$ : (a : A) (f : Pi A B), f = $\lambda$ (x : A), f x) \\
             & ($\Sigma$-form$_1$ : U) \\
             & ($\Sigma$-ctor$_1$ : $\Pi$ (a : A) (b : B a) , Sigma A B) \\
             & ($\Sigma$-elim$_1$ : $\Pi$ (p : Sigma A B), A) \\
             & ($\Sigma$-elim$_2$ : $\Pi$ (p : Sigma A B), B (pr$_1$ A B p)) \\
             & ($\Sigma$-comp$_1$ : $\Pi$ (a : A) (b: B a), a = $\Sigma$-elim$_1$ ($\Sigma$-ctor$_1$ a b)) \\
             & ($\Sigma$-comp$_2$ : $\Pi$ (a : A) (b: B a), b = $\Sigma$-elim$_2$ (a, b)) \\
             & ($\Sigma$-comp$_3$ : $\Pi$ (p : Sigma A B), p = (pr$_1$ A B p, pr$_2$ A B p)), \textbf{1}
\end{tabular}
\end{table}

\newpage
\subsection{Path Equality}

The fundamental development of equality inside MLTT provers led us to the
notion of $\infty$-groupoid as spaces. In this way Path identity type appeared
in the core of type checker along with De Morgan algebra on built-in interval type.

\begin{table}[ht!]
\tabstyle
\begin{tabular}{rl}
     \textbf{def} & homotopy (lang : U) : \textbf{U} := \textbf{inductive} \\
               \{ & PathP: lang $\rightarrow$ lang $\rightarrow$ lang \\
  |& plam: name $\rightarrow$ lang $\rightarrow$ lang \\
  |& papp: lang $\rightarrow$ lang \\
  |& I \\
  |& zero \\
  |& one \\
  |& meet: lang $\rightarrow$ lang \\
  |& join: lang $\rightarrow$ lang \\
  |& neg: lang \\
  |& system: lang \\
  |& Partial: lang \\
  |& transp: lang $\rightarrow$ lang \\
  |& hcomp: lang \\
  |& Sub: lang \\
  |& inc: lang \\
  |& ouc: lang \\
 \}& \\
\end{tabular}
\end{table}

\begin{definition}[Cubical Presheaf $\mathbb{I}$]
The identity types modeled with another presheaf, the presheaf on Lawvere
category of distributive lattices (theory of De Morgan algebras) denoted
with $\Box$ — $\textbf{I} : \Box^{op} \rightarrow \textrm{Set}$.
\end{definition}

\begin{definition}[Properties of $\textbf{I}$] The presheaf $\textbf{I}$:
i) has to distinct global elements $0$ and $1$ (B$_1$);
ii) $\textbf{I}$(I) has a decidable equality for each $I$ (B$_2$);
iii) $\textbf{I}$ is tiny so the path functor $X \mapsto X^\textbf{I}$ has right adjoint (B$_3$).;
iv) $\textbf{I}$ has meet and join (connections).
\end{definition}

$\textbf{Interval\ Pretypes}$. While having pretypes universe V with interval and
associated De Morgan algebra ($\wedge$, $\vee$, -, 0, 1, I) is enough to
perform DNF normalization and proving some basic statements about path, including:
contractability of singletons, homotopy transport, congruence, functional
extensionality; it is not enough for proving $\beta$ rule for Path type or path composition.
\\
\\
\indent $\textbf{Generalized\ Transport}$. Generalized transport transp adresses
first problem of deriving the computational $\beta$ rule for Path types:

\begin{table}[ht!]
\tabstyle
\begin{tabular}{rl}
\textbf{theorem} & Path$_\beta$ (A : U) (a : A) (C : D A) (d: C a a (refl A a)) \\
               : & Equ (C a a (refl A a)) d (J A a C d a (refl A a)) \\
              := & $\lambda$ (A : U), $\lambda$ (a : A), \\
                 & $\lambda$ (C : $\Pi$ (x : A), $\Pi$ (y : A), PathP (<\_> A) x y $\rightarrow$ U), \\
                 & $\lambda$ (d : C a a (<\_> a)), <j> transp (<\_> C a a (<\_> a)) -j d \\
\end{tabular}
\end{table}

Transport is defined on fibrant types (only) and type checker should cover all the cases
Note that transp$^i$ (Path$^j$ A v w) $\varphi$ u$_0$ case is relying on comp
operation which depends on hcomp primitive. Here is given the first part of Simon Huber equations \cite{Huber} for \textbf{transp}:

\begin{table}[ht!]
\tabstyle
\begin{tabular}{l}
transp$^{i}$ N $\varphi$ u$_0$ = u$_0$ \\
transp$^{i}$ U $\varphi$ A = A \\
transp$^{i}$ ($\Pi$ (x : A), B) $\varphi$ u$_0$ v = transp$^i$ B(x/w) $\varphi$ (u$_0$ w(i/0)) \\
transp$^{i}$ ($\Sigma$ (x : A), B) $\varphi$ u$_0$ = (transp$^i$ A $\varphi$ (u$_0$.1),transp$^i$ B(x/v) $\varphi$ (u$_0$.2)) \\
transp$^{i}$ (Path$^j$ v w) $\varphi$ u$_0$ = <j> comp$^i$ A [$\phi$ u$_0$ j, (j=0) $\mapsto$ v, (j=1) $\mapsto$ w] (u$_0$ j) \\
transp$^{i}$ (Glue [$\varphi$ $\mapsto$ (T,w)] A) $\psi$ u$_0$ = glue [$\phi$(i/1) $\mapsto$ t'$_1$] a'$_1$ : B(i/1) \\
\end{tabular}
\end{table}

\newpage
$\textbf{Partial\ Elements}$. In order to explicitly define hcomp we need to specify
n-cubes where some faces are missing. Partial primitives isOne, 1=1 and UIP on pretypes
are derivable in Anders due to landing strict equality Id in V universe. The idea is
that (Partial A r) is the type of cubes in A that are only defined when IsOne r holds.
(Partial A r) is a special version of the function space IsOne r → A with a more
extensional equality: two of its elements are considered judgmentally equal if
they represent the same subcube of A. They are equal whenever they reduce to
equal terms for all the possible assignment of variables that make r equal to 1.

\begin{table}[ht!]
\tabstyle
\begin{tabular}{l}
\textbf{def} Partial' (A : U) (i : I) := Partial A i \\
\textbf{def} isOne : I -> V := Id I 1 \\
\textbf{def} 1=>1 : isOne 1 := ref 1 \\
\textbf{def} UIP (A : V) (a b : A) (p q : Id A a b) : Id (Id A a b) p q := ref p \\
\end{tabular}
\end{table}

$\textbf{Cubical\ Subtypes}$. For (A : U) (i : I) (Partial A i) we can define
subtype A [ i $\mapsto$ u ]. A term of this type is a term of type A that is
definitionally equal to u when (IsOne i) is satisfied. We have forth and back
fusion rules ouc (inc v) = v and inc (outc v) = v. Moreover, ouc v will reduce to u 1=1 when i=1.

\begin{table}[ht!]
\tabstyle
\begin{tabular}{l}
\textbf{def} sub' (A : U) (i : I) (u : Partial A i) : V := A [i $\mapsto$ u ] \\
\textbf{def} inc' (A : U) (i : I) (a : A) : A [i $\mapsto$ [(i = 1) $\rightarrow$ a]] := inc A i a \\
\textbf{def} ouc' (A : U) (i : I) (u : Partial A i) (a : A [i $\mapsto$ u]) : A := ouc a \\
\end{tabular}
\end{table}

$\textbf{Homogeneous\ Composition}$. hcomp is the answer to second problem: with hcomp and transp
one can express path composition, groupoid, category of groupoids (groupoid interpretation and
internalization in type theory). One of the main roles of homogeneous composition is to be a
carrier in [higher] inductive type constructors for calculating of homotopy colimits and direct
encoding of CW-complexes. Here is given the second part of Simon Huber equations \cite{Huber} for $\textbf{hcomp}$:

\begin{table}[ht!]
\tabstyle
\begin{tabular}{l}
hcomp$^i$ N [$\phi$ $\mapsto$ 0] 0 = 0 \\
hcomp$^i$ N [$\phi$ $\mapsto$ S u] (S u$_0$) = S (hcomp$^i$ N [$\phi$ $\mapsto$ u] u$_0$) \\
hcomp$^i$ U [$\phi$ $\mapsto$ E] A = Glue [$\phi$ $\mapsto$ (E(i/1), equiv$^i$ E(i/1-i))] A \\
hcomp$^i$ ($\Pi$ (x : A), B) [$\phi$ $\mapsto$ u] u$_0$ v = hcomp$^i$ B(x/v) [$\phi$ $\mapsto$ u v] (u$_0$ v) \\
hcomp$^i$ ($\Sigma$ (x : A), B) [$\phi$ $\mapsto$ u] u$_0$ = (v(i/1), comp$^i$ B(x/v) [$\phi$ $\mapsto$ u.2] u$_0$.2) \\
hcomp$^i$ (Path$^j$ A v w) [$\phi$ $\mapsto$ u] u$_0$ = <j> hcomp$^i$ A[$\phi$ $\mapsto$ u j, (j=0) $\mapsto$ v, (j=1) $\mapsto$ w] (u$_0$ j) \\
hcomp$^i$ (Glue [$\phi$ $\mapsto$ (T,w)] A) [$\psi$ $\mapsto$ u] u$_0$ = glue [$\phi$ $\mapsto$ u(i/1)] (unglue u(i/1)) \\
\end{tabular}
\end{table}

\subsection{Strict Equality}

To avoid conflicts with path equalities which live in fibrant universes strict
equalities live in pretypes universes.

\begin{table}[ht!]
\tabstyle
\begin{tabular}{rl}
     \textbf{def} & strict (lang : U) : \textbf{U} := \textbf{inductive} \\\
 \{ & Id: name $\rightarrow$ lang \\
  |& ref: lang $\rightarrow$ lang \\
  |& idJ: lang $\rightarrow$ lang $\rightarrow$ lang \\
                         \} & \\
\end{tabular}
\end{table}

We use strict equality in HTS for pretypes and partial elements which live in V.
The presheaf configuration with Pi, Sigma and Id is called \textbf{MLTT-75} \cite{MLTT75}. The presheaf
configuration with Pi, Sigma, Id and Path is called \textbf{HTS} (Homotopy Type System).

\newpage
\subsection{Glue Types}

The main purpose of Glue types is to construct a cube where some faces have
been replaced by equivalent types. This is analogous to how hcomp lets us
replace some faces of a cube by composing it with other cubes, but for Glue
types you can compose with equivalences instead of paths. This implies the
univalence principle and it is what lets us transport along paths built
out of equivalences.

\begin{table}[ht!]
\tabstyle
\begin{tabular}{rl}
     \textbf{def} & glue (lang : U) : \textbf{U} := \textbf{inductive} \\
 \{ & Glue: lang $\rightarrow$ lang $\rightarrow$ lang \\
  |& glue: lang $\rightarrow$ lang \\
  |& unglue: lang $\rightarrow$ lang \\
                         \} & \\
\end{tabular}
\end{table}

Basic Fibrational HoTT core by Pelayo, Warren, and Voevodsky (2012).

\begin{table}[ht!]
\tabstyle
\begin{tabular}{l}
\textbf{def} fiber (A B : U) (f: A → B) (y : B): U := $\Sigma$ (x : A), Path B (f x) y\\
\textbf{def} isEquiv (A B : U) (f : A → B) : U := $\Pi$ (y : B), isContr (fiber A B f y) \\
\textbf{def} equiv (A B : U) : U := $\Sigma$ (f : A → B), isEquiv A B f \\
\textbf{def} contrSingl (A : U) (a b : A) (p : Path A a b) \\
\ \ \ \ : Path ($\Sigma$ (x : A), Path A a x) (a,<i>a) (b,p) := <i> (p @ i, <j> p @ i $\vee$ j) \\
\textbf{def} idIsEquiv (A : U) : isEquiv A A (id A) := \\
\ \ \ \ $\lambda$ (a : A), ((a,<i>a), $\lambda$ (z : fiber A A (id A) a), contrSingl A a z.1 z.2) \\
\textbf{def} idEquiv (A : U) : equiv A A := (id A, isContrSingl A) \\
\end{tabular}
\end{table}

The notion of Univalence was discovered by Vladimir Voevodsky
as forth and back transport between fibrational equivalence
as contractability of fibers and homotopical multi-dimentional
heterogeneous path equality. The Equiv → Path type is called Univalence type,
where univalence intro is obtained by Glue type and elim (Path → Equiv) is
obtained by sigma transport from constant map.

\begin{table}[ht!]
\tabstyle
\begin{tabular}{l}
\textbf{def} univ-formation (A B : U) := equiv A B → PathP (<i> U) A B \\
\textbf{def} univ-intro (A B : U) : univ-formation A B := $\lambda$ (e : equiv A B), \\
\ \ \ \ <i> Glue B ($\partial$ i) [(i = 0) → (A, e), (i = 1) → (B, idEquiv B)] \\
\textbf{def} univ-elim (A B : U) (p : PathP (<i> U) A B) \\
\ \ : equiv A B := transp (<i> equiv A (p @ i)) 0 (idEquiv A) \\
\textbf{def} univ-computation (A B : U) (p : PathP (<i> U) A B) \\
\ \ : PathP (<i> PathP (<i> U) A B) (univ-intro A B (univ-elim A B p)) p \\
\ \ := <j i> Glue B (j $\vee$ $\partial$ i) \\
\ \ \ \ \ [(i = 0) → (A, univ-elim A B p), (i = 1) → (B, idEquiv B), \\
\ \ \ \ \ \ (j = 1) → (p @ i, univ-elim (p @ i) B (<k> p @ (i $\vee$ k)))] \\
\end{tabular}
\end{table}

Similar to Fibrational Equivalence the notion of Retract/Section based Isomorphism could be introduced
as forth-back transport between isomorphism and path equality. This notion is somehow cannonical to all
cubical systems and is called Unimorphism here.

\begin{table}[ht!]
\tabstyle
\begin{tabular}{l}
\textbf{def} iso-Form (A B: U) : U$_1$ := iso A B -> PathP (<i>U) A B \\
\textbf{def} iso-Intro (A B: U) : iso-Form A B := \\
\ \ \ $\lambda$ (x : iso A B), isoPath A B x.f x.g x.s x.t \\
\textbf{def} iso-Elim (A B : U) : PathP (<i> U) A B -> iso A B \\
\ := $\lambda$ (p : PathP (<i> U) A B), \\
\ \ \ ( coerce A B p, coerce B A (<i> p @ -i), \\
\ \ \ \ \ trans$^{-1}$-trans A B p, $\lambda$ (a : A), <k> trans-trans$^{-1}$ A B p a @-k, $\star$) \\
\end{tabular}
\end{table}

Orton-Pitts basis for univalence computability (2017):

\begin{table}[ht!]
\tabstyle
\begin{tabular}{l}
\textbf{def} ua (A B : U) (p : equiv A B) : PathP (<i> U) A B := univ-intro A B p \\
\textbf{def} ua$-\beta$ (A B : U) (e : equiv A B) : Path (A → B) (trans A B (ua A B e)) e.1 \\
 := <i> $\lambda$ (x : A), (idfun=idfun'' B @ -i) \\
\ \ \ \ ( (idfun=idfun'' B @ -i) ((idfun=idfun' B @ -i) (e.1 x)) ) \\
\end{tabular}
\end{table}

\newpage
\subsection{de Rham Stack}

Stack de Rham or Infinitezemal Shape Modality is a basic primitive for proving theorems
from synthetic differential geometry. This type-theoretical framework was developed
for the first time by Felix Cherubini under the guidance of Urs Schreiber. The Anders
prover implements the computational semantics of the de Rham stack.

\begin{table}[ht!]
\tabstyle
\begin{tabular}{l}
\textbf{def} $\iota$ (A : U) (a : A) : $\Im$ A := $\Im-$unit a \\
\textbf{def} $\mu$ (A : U) (a : $\Im$ ($\Im$ A)) := $\Im-$join a \\
\\
\textbf{def} is-coreduced (A : U) : U := isEquiv A ($\Im$ A) ($\iota$ A) \\
\textbf{def} $\Im$-coreduced (A : U) : is-coreduced ($\Im$ A) := isoToEquiv \\
\ \ \ \ ($\Im$ A) ($\Im$ ($\Im$ A)) ($\iota$ ($\Im$ A)) ($\mu$ A) ($\lambda$ (x : $\Im$ ($\Im$ A)), <i>x) ($\lambda$ (y : $\Im$ A),<i>y) \\
\textbf{def} ind$-\Im\beta$ (A : U) (B : $\Im$ A → U) (f : $\Pi$ (a : A), $\Im$ (B ($\iota$ A a))) (a : A) \\
\ \ : Path ($\Im$ (B ($\iota$ A a))) (ind$-\Im$ A B f ($\iota$ A a)) (f a) := <i> f a \\
\textbf{def} ind$-\Im-$const (A B : U) (b : $\Im$ B) (x : $\Im$ A) \\
\ \ : Path ($\Im$ B) (ind$-\Im$ A ($\lambda$ (i : $\Im$ A), B) ($\lambda$ (i : A), b) x) b := <i> b \\
\end{tabular}
\end{table}

Coreduced induction and its $\beta-$quation.

\begin{table}[ht!]
\tabstyle
\begin{tabular}{l}
\textbf{def} $\Im-$ind (A : U) (B : $\Im$ A → U) (c : $\Pi$ (a : $\Im$ A), \\
\ \ \ \ is-coreduced (B a)) (f : $\Pi$ (a : A), B ($\iota$ A a)) (a : $\Im$ A) : B a \\
\ \ \ \ := (c a (ind$-\Im$ A B ($\lambda$ (x : A), $\iota$ (B ($\iota$ A x)) (f x)) a)).1.1 \\
\textbf{def} $\Im-$ind$\beta$ (A : U) (B : $\Im$ A → U) (c : $\Pi$ (a : $\Im$ A), \\
\ \ \ \ is-coreduced (B a)) (f : $\Pi$ (a : A), B ($\iota$ A a)) (a : A) \\
\ \ : Path (B ($\iota$ A a)) (f a) (($\Im-$ind A B c f) ($\iota$ A a)) \\
\ := <i> sec-equiv (B ($\iota$ A a)) ($\Im$ (B ($\iota$ A a))) ($\iota$ (B ($\iota$ A a)), c ($\iota$ A a)) (f a) @-i \\
\end{tabular}
\end{table}

Geometric Modal HoTT Framework: Infinitesimal Proximity, Formal Disk, Formal Disk Bundle, Differential.

\begin{table}[ht!]
\tabstyle
\begin{tabular}{l}
\textbf{def} $\sim$ (X : U) (a x' : X) : U := Path ($\Im$ X) ($\iota$ X a) ($\iota$ X x') \\
\textbf{def} $\mathbb{D}$ (X : U) (a : X) : U := $\Sigma$ (x' : X), $\sim$ X a x' \\
\textbf{def} inf-prox-ap (X Y : U) (f : X → Y) (x x' : X) (p : $\sim$ X x x') \\
\ \ : $\sim$ Y (f x) (f x') := <i> $\Im-$app X Y f (p @ i) \\
\textbf{def} T$^\infty$ (A : U) : U := $\Sigma$ (a : A), $\mathbb{D}$ A a \\
\textbf{def} inf-prox-ap (X Y : U) (f : X → Y) (x x' : X) (p : $\sim$ X x x') \\
\ \ : $\sim$ Y (f x) (f x') := <i> $\Im-$app X Y f (p @ i) \\
\textbf{def} d (X Y : U) (f : X → Y) (x : X) ($\varepsilon$ : $\mathbb{D}$ X x) \\
\ \ : $\mathbb{D}$ Y (f x) := (f $\varepsilon$.1, inf-prox-ap X Y f x $\varepsilon$.1 $\varepsilon$.2) \\
\textbf{def} T$^\infty$-map (X Y : U) (f : X → Y) ($\tau$ : T$^\infty$ X) : T$^\infty$ Y := (f $\tau$.1, d X Y f $\tau$.1 $\tau$.2) \\
\end{tabular}
\end{table}

\subsection{Flat Modality}

The Flat modality $\flat$ (also known as the discrete modality) is a core primitive of Anders,
providing a way to talk about discrete types. In a cohesive type theory, $\flat A$
is the type $A$ equipped with the discrete topology. The Anders prover implements
computational semantics for the Flat modality including its interaction with cubical primitives.

\begin{table}[ht!]
\tabstyle
\begin{tabular}{l}
\textbf{def} Flat (A : U) := $\flat$ A \\
\textbf{def} FlatUnit (A : U) (a : A) := $\flat$-unit a \\
\textbf{def} FlatCounit (A : U) (x : $\flat$ A) := $\flat$-counit x \\
\\
\textbf{def} isDiscrete (A : U) : U := isEquiv A ($\flat$ A) (FlatUnit A) \\
\textbf{def} $\flat$-$\eta$ (A : U) (x : $\flat$ A) : Path ($\flat$ A) x ($\flat$-unit ($\flat$-counit x)) := <p> x \\
\textbf{def} FlatInd (A : U) (B : $\flat$ A $\rightarrow$ U) := ind-$\flat$ A B \\
\end{tabular}
\end{table}

The Flat modality satisfies several computational reduction rules for transport and
homogeneous composition, which ensure that $\flat$ behaves correctly with respect to the
cubical structure of the type theory.

\begin{table}[ht!]
\tabstyle
\begin{tabular}{l}
transp$^{i}$ ($\flat$ A) $\varphi$ ($\flat$-unit a) = $\flat$-unit (transp$^i$ A $\varphi$ a) \\
hcomp$^{i}$ ($\flat$ A) [$\phi$ $\mapsto$ $\flat$-unit u] ($\flat$-unit u$_0$) = $\flat$-unit (hcomp$^i$ A [$\phi$ $\mapsto$ u] u$_0$) \\
\end{tabular}
\end{table}

\begin{definition}[Discreteness]
A type $A$ is called discrete if the unit map $\flat$-unit : $A \rightarrow \flat A$ is an equivalence.
This corresponds to the intuition that a type is discrete if its points have no non-trivial paths between them
that aren't captured by the discrete topology.
\end{definition}

\subsection{Inductive Types}

The foundational inductive types in \anders\ are based on Martin-L\"{o}f's intuitionistic type theory \cite{ML84}. Instead of a generic inductive definition scheme, \anders\ provides well-founded trees (W-types) as a primitive construction from which other data types can be derived. Basic data types like List, Fin, and Vec are implemented as W-types in the core library.

\subsubsection*{Empty Type}

The empty type $\mathbf{0}$ (also written as \texttt{empty}) is the inductive type with no constructors. It represents logical falsehood. Its induction principle is the principle of explosion (\textit{ex falso quodlibet}).

\begin{table}[ht!]
\tabstyle
\begin{tabular}{l}
\textbf{def} Empty : U := \textbf{0} \\
\textbf{def} Empty-elim (C : Empty $\rightarrow$ U) (z : Empty) : C z := \textbf{ind-empty} C z
\end{tabular}
\end{table}

\subsubsection*{Unit Type}

The unit type $\mathbf{1}$ (also written as \texttt{unit}) is the inductive type with a single constructor ★. It represents logical truth.

\begin{table}[ht!]
\tabstyle
\begin{tabular}{l}
\textbf{def} Unit : U := \textbf{1} \\
\textbf{def} star : Unit := ★ \\
\\
\textbf{def} Unit-ind (C : Unit $\rightarrow$ U) (u : C star) (z : Unit) : C z := \textbf{ind-unit} C u z
\end{tabular}
\end{table}

\subsubsection*{Boolean Type}

The boolean type $\mathbf{2}$ (also written as \texttt{bool}) is the inductive type with two constructors $0_2$ and $1_2$ (or \texttt{false} and \texttt{true}). It is used for case analysis and decision procedures.

\begin{table}[ht!]
\tabstyle
\begin{tabular}{l}
\textbf{def} Bool : U := \textbf{2} \\
\textbf{def} false : Bool := $0_2$ \\
\textbf{def} true : Bool := $1_2$ \\
\\
\textbf{def} Bool-ind (C : Bool $\rightarrow$ U) (f : C false) (t : C true) (z : Bool) \\
\ \ : C z := \textbf{ind-bool} C f t z
\end{tabular}
\end{table}

\newpage
\subsubsection*{Natural Numbers}

The type of natural numbers $\mathbb{N}$ is a primitive type in \anders, providing efficient normalization for arithmetic operations. Its elimination rule is provided by the \texttt{ind-Nat} primitive.

\begin{table}[ht!]
\tabstyle
\begin{tabular}{l}
\textbf{def} Nat : U := \textbf{Nat} \\
\textbf{def} Zero : Nat := \textbf{0} \\
\textbf{def} Succ (n : Nat) : Nat := \textbf{S} n \\
\\
\textbf{def} Nat-ind (C : Nat $\rightarrow$ U) (z : C Zero) (s : $\Pi$ (n : Nat), C n $\rightarrow$ C (Succ n)) \\
\ \ (n : Nat) : C n := \textbf{ind-Nat} C z s n
\end{tabular}
\end{table}

\subsubsection*{W-Types}

Well-founded trees (W-types) are the building blocks for all other inductive types in \anders. A W-type $\mathsf{W}(x : A), B(x)$ is specified by a type $A$ of shapes and a family $B : A \to U$ of positions for each shape.

\begin{table}[ht!]
\tabstyle
\begin{tabular}{l}
\textbf{def} W (A : U) (B : A $\rightarrow$ U) : U := \textbf{W} A B \\
\textbf{def} sup (A : U) (B : A $\rightarrow$ U) (a : A) (f : B a $\rightarrow$ W A B) \\
\ \ : W A B := \textbf{sup} a f \\
\\
\textbf{def} W-ind (A : U) (B : A $\rightarrow$ U) (C : W A B $\rightarrow$ U) \\
\ \ (g : $\Pi$ (a : A) (f : B a $\rightarrow$ W A B), ($\Pi$ (b : B a), C (f b)) $\rightarrow$ C (sup A B a f)) \\
\ \ (w : W A B) : C w := \textbf{ind-W} A B C g w
\end{tabular}
\end{table}

\newpage
\subsection{Higher Inductive Types}

Higher inductive types (HITs) generalize inductive types by allowing constructors that target path types, enabling the representation of cell complexes and other quotient-like structures. \anders\ provides two primitive HIT constructions: coequalizers and hub-and-spoke discs. These primitives, combined with W-types, are sufficient to represent a wide range of HITs, including spheres $S^n$, suspensions, and truncations, following the Three-HIT theorem approach \cite{BauerWeide}.

\subsubsection*{Coequalizers}

The coequalizer $\mathsf{coequ}(f, g)$ of two parallel maps $f, g : A \to B$ is a higher inductive type that identifies the images of $f$ and $g$ in $B$. That is, for every $a : A$, we have a path $f(a) = g(a)$ in $\mathsf{coequ}(f, g)$. This allows for computable representations of the circle $S^1$, the torus, and general colimits.

\begin{table}[ht!]
\tabstyle
\begin{tabular}{l}
\textbf{def} coeq (A B : U) (f g : A $\rightarrow$ B) : U := \textbf{coequ} A B f g \\
\textbf{def} iota (A B : U) (f g : A $\rightarrow$ B) (x : B) : \textbf{coequ} A B f g := $\iota_2$ A B f g x \\
\textbf{def} resp (A B : U) (f g : A $\rightarrow$ B) (x : A) \\
\ \ : Path (\textbf{coequ} A B f g) ($\iota_2$ A B f g (f x)) ($\iota_2$ A B f g (g x)) \\
\ := <i> \textbf{resp} A B f g x @ i \\
\\
\textbf{def} coequ-ind' (A B : U) (f g : A $\rightarrow$ B) (X : \textbf{coequ} A B f g $\rightarrow$ U) \\
\ \ (i : $\Pi$ (b : B), X (iota A B f g b)) \\
\ \ ($\rho$ : $\Pi$ (x : A), PathP (<i> X (resp A B f g x @ i)) (i (f x)) (i (g x))) \\
\ \ (z : \textbf{coequ} A B f g) : X z \\
\ := coequ-ind A B f g X i $\rho$ z

\end{tabular}
\end{table}

\subsubsection*{Hub-and-Spoke Discs}

The hub-and-spoke construction $\mathsf{disc}(S, A)$ attaches a "disc" to a type $A$ for every map $f : S \to \mathsf{disc}(S, A)$. This is achieved by adding a "hub" point for each such $f$ and a "spoke" path connecting the hub to $f(s)$ for every $s : S$. This construction is particularly useful for defining $n$-truncations and other cell complexes where a point must be connected to an entire family of points.

\begin{table}[ht!]
\tabstyle
\begin{tabular}{l}
\textbf{def} hs (S A : U) : U := \textbf{disc} S A \\
\textbf{def} center (S A : U) (a : A) : hs S A := \textbf{base} S A a \\
\textbf{def} hub' (S A : U) (f : S $\rightarrow$ hs S A) : hs S A := \textbf{hub} S A f \\
\textbf{def} spoke' (S A : U) (f : S $\rightarrow$ hs S A) (s : S) \\
\ \ : PathP (<i> hs S A) (hub' S A f) (f s) := \textbf{spoke} S A f s \\
\\
\textbf{def} hs-ind (S A : U) (X : hs S A $\rightarrow$ U) \\
\ \ (nCenter : $\Pi$ (a : A), X (center S A a)) \\
\ \ (nHub : $\Pi$ (f : S $\rightarrow$ hs S A) (nF : $\Pi$ (s : S), X (f s)), X (hub' S A f)) \\
\ \ (nSpoke : $\Pi$ (f : S $\rightarrow$ hs S A) (nF : $\Pi$ (s : S), X (f s)) (s : S), \\
\ \ \ \ PathP (<i> X (spoke' S A f s @ i)) (nHub f nF) (nF s)) \\
\ \ (z : hs S A) : X z := \textbf{disc-ind} S A X nCenter nHub nSpoke z
\end{tabular}
\end{table}


% \subsection{Simplicial Types}

% Modification of Anders with Simplicial types and Hopf Fibrations built intro the core of type checker
% is called \textbf{Dan} with following recursive syntax (having $f$ as Simplecies and $coh$ as Path-coherence functions):

% \begin{table}[ht!]
% \tabstyle
% \begin{tabular}{l}
%    \textbf{simplex} n [v$_0$ .. v$_n$] \{ f$_0$, f$_1$, ..., f$_n$ | coh i$_1$ i$_2$ ... i$_n$ \} : Simplex
% \end{tabular}
% \end{table}
% 
% and instantiation example:

% \begin{table}[ht!]
% \tabstyle
% \begin{tabular}{l}
% 
% \textbf{def} s$_\infty$ : Simplicial \\
% \ := $\Pi$ (v e : Simplex), \\
% \ \ \ \ $\delta_{10}$ = v, $\delta_{11}$ = v, s$_0$ < v, \\
% \ \ \ \ $\delta_{20}$ = e $\circ$ e, s$_{10}$ < $\delta_{20}$ \\
% \ \ \ \ $\vdash$ $\infty$ (v, e, $\delta_{20}$ | $\delta_{10}$ $\delta_{11}$, s$_0$, $\delta_{20}$, s$_{10}$)
% \end{tabular}
% \end{table}

\newpage
\section{Properties}

Soundness and completeness link syntax to semantics.
Canonicity, normalization, and totality ensure computational adequacy.
Consistency and decidability guarantee logical and practical usability.
Conservativity and initiality support extensibility and universality.

\subsection{Soundness and Completeness}

Soundness is proven via cubical sets \cite{CCHM, Awodey12, Coquand18}.

\subsection{Canonicity, Normalization and Totality}

Canonicity and normalization hold constructively \cite{Huber17, Streicher91}.

\subsection{Consistency and Decidability}

Consistency follows from the model \cite{Bezem14}.
Decidability is achieved for type checking \cite{Coquand18}.

\subsection{Conservativity and Initiality}

Conservativity and initiality is discussed bu Shulman\cite{Hofmann97, Shulman15}.
Initiality is implicit in the syntactic construction \cite{Awodey12}.

\section{Conclusion}

This paper presents Anders, a proof assistant that reimplements
cubicaltt within a Modal Homotopy Type System framework, based
on MLTT-80 and CCHM/CHM. It integrates HTS strict equality,
infinitesimal modalities, and primitives like suspensions or quotients,
with the extension adding simplicial types and Hopf fibrations.
Anders offers an efficient, idiomatic system — compiling in under
one second — using a syntax of Lean and semantics of cubicaltt and Cubical Agda.
As a practical refinement of cubicaltt, Anders serves as an accessible
tool for homotopy type theory, with potential for incremental
enhancements like a tactic language.

\appendix
\newpage
\section{Annex A. Simon Huber Canonical Equations}
\label{sec:appendix-b}

\begin{table}[ht!]
\tabstyle
\begin{tabular}{l}
transpⁱ N φ u₀ = u₀ \\
transpⁱ U φ A = A \\
transpⁱ (Π (x : A), B) φ u₀ v = transpⁱ B(x/w) φ (u₀ w(i/0)), w = tFill⁻ⁱ A φ v, v : A(i/1) \\
transpⁱ (Σ (x : A), B) φ u₀ = (transpⁱ A φ (u₀.1), transpⁱ B(x/v) φ(u₀.2)), v = tFillⁱ A φ u₀.1 \\
transpⁱ (Pathʲ A v w) φ u₀ = 〈j〉\\
\ compⁱ A [φ ↦ u₀ j, (j=0) ↦ v, (j=1) ↦ w] (u₀ j) \\
transpⁱ G ψ u₀ = glue [φ(i/1) ↦ t′₁] a′₁ : G(i/1), \\
\ G = Glue [φ ↦ (T,w)] A \\
transpⁱ (W (x : A), B) φ (sup a f) = \\
\ sup (transpⁱ A φ a) (transpⁱ (B(v) → W) φ f), v = tFillⁱ A φ a \\
transpⁱ (Coequ A B f g) φ (ι₂ b) = ι₂ (transpⁱ B φ b) \\
transpⁱ (Coequ A B f g) φ (resp a j) = resp (transpⁱ A φ a) j \\
transpⁱ (Disc S A) φ (base a) = base (transpⁱ A φ a) \\
transpⁱ (Disc S A) φ (hub f) = hub (transpⁱ (S → Disc S A) φ f) \\
transpⁱ (Disc S A) φ (spoke f y j) = spoke (transpⁱ (S → Disc S A) φ f) (transpⁱ S φ y) j \\
transpⁱ (ℑ A) φ (ℑ-unit a) = ℑ-unit (transpⁱ A φ a) \\
transpⁱ (♭ A) φ (♭-unit a) = ♭-unit (transpⁱ A φ a) \\
transpⁱ (♯ A) φ (♯-unit a) = ♯-unit (transpⁱ A φ a) \\
transpⁱ (tw A) φ (tw-unit a) = tw-unit (transpⁱ A φ a) \\
transpⁱ (Π (x :\textsubscript{μ} A), B) φ u₀ v = transpⁱ B[x/w] φ (u₀ w(i/0)) \\
transp⁻ⁱ A φ u = (transpⁱ A(i/1−i) φ u)(i/1−i) : A(i/0) \\
tFillⁱ A φ u₀ = transpʲ A(i/i∧j) (φ∨(i=0)) u₀ : A \\
\\
hcompⁱ N [φ ↦ 0] 0 = 0 \\
hcompⁱ N [φ ↦ S u] (S u₀) = S (hcompⁱ N [φ ↦ u] u₀) \\
hcompⁱ U [φ ↦ E] A = Glue [φ ↦ (E(i/1), equivⁱ E(i/1−i))] A \\
hcompⁱ (Π (x : A), B) [φ ↦ u] u₀ v = hcompⁱ B(x/v) [φ ↦ u v] (u₀ v) \\
hcompⁱ (Σ (x : A), B) [φ ↦ u] u₀ = (v(i/1), compⁱ B(x/v) [φ ↦ u.2] u₀.2), \\
\ v = hFillⁱ A [φ ↦ u.1] u₀.1 \\
hcompⁱ (Pathʲ A v w) [φ ↦ u] u₀ = 〈j〉\\
\ hcompⁱ A [ φ ↦ u j, (j = 0) ↦ v, (j = 1) ↦ w ] (u₀ j) \\
hcompⁱ G [ψ ↦ u] u₀ = glue [φ ↦ t₁] a₁ : G, G = Glue [φ ↦ (T,w)] A, \\
t₁ = u(i/1) : T, a₁ = unglue u(i/1) : A, glue [φ ↦ t₁] a₁ = t₁ : T \\
hcompⁱ (W (x : A), B) [φ ↦ sup a f] (sup a₀ f₀) \\
\ = sup (hcompⁱ A [φ ↦ a] a₀) (hcompⁱ (B(v) → W) [φ ↦ f] f₀), v = hFillⁱ A [φ ↦ a] a₀ \\
hcompⁱ (Coequ A B f g) [φ ↦ ι₂ u] (ι₂ u₀) = ι₂ (hcompⁱ B [φ ↦ u] u₀) \\
hcompⁱ (Coequ A B f g) [φ ↦ resp a j] (resp a₀ j) = resp (hcompⁱ A [φ ↦ a] a₀) j \\
hcompⁱ (Disc S A) [φ ↦ base a] (base a₀) = base (hcompⁱ A [φ ↦ a] a₀) \\
hcompⁱ (Disc S A) [φ ↦ hub f] (hub f₀) = hub (hcompⁱ (S → Disc S A) [φ ↦ f] f₀) \\
hcompⁱ (Disc S A) [φ ↦ spoke f y j] (spoke f₀ y₀ j) \\
\ = spoke (hcompⁱ (S → Disc S A) [φ ↦ f] f₀) (hcompⁱ S [φ ↦ y] y₀) j \\
hcompⁱ (ℑ A) [φ ↦ ℑ-unit u] (ℑ-unit u₀) = ℑ-unit (hcompⁱ A [φ ↦ u] u₀) \\
hcompⁱ (♭ A) [φ ↦ ♭-unit u] (♭-unit u₀) = ♭-unit (hcompⁱ A [φ ↦ u] u₀) \\
hcompⁱ (♯ A) [φ ↦ ♯-unit u] (♯-unit u₀) = ♯-unit (hcompⁱ A [φ ↦ u] u₀) \\
hcompⁱ (tw A) [φ ↦ tw-unit u] (tw-unit u₀) = tw-unit (hcompⁱ A [φ ↦ u] u₀) \\
hcompⁱ (Π (x :\textsubscript{μ} A), B) [φ ↦ u] u₀ v = hcompⁱ B[x/v] [φ ↦ u v] (u₀ v) \\
hFillⁱ A [φ ↦ u] u₀ = hcompʲ A [φ ↦ u(i/i∧j), (i=0) ↦ u₀] u₀ : A \\
\end{tabular}
\end{table}

\newpage
\section{Annex B. HoTT Mathematics}
\label{sec:hott-applications}

Mathematics in Homotopy Type Theory has seen substantial formalization efforts across multiple proof assistants (including Anders, Cubical Agda, Lean, Coq, etc.). Below is a (non-exhaustive) list of notable results that have been mechanized.

\subsection{Formalized Results}

\begin{itemize}
\item Symmetry: A book project about group theory from the univalent point of view~\cite{SymmetryBook}.
\item Modalities and reflective subuniverses~\cite{RijkeModalities}.
\item Higher Groups ($k$-symmetric $n$-groups)~\cite{BuchholtzHigherGroups}.
\item $\lVert \Sigma A \rVert_2$ is a 1-type for all sets $A$~\cite{KrausFreeHigherGroups}.
\item Localization and $L$-separated types~\cite{ChristensenLocalization}.
\item The James construction and $\pi_4(S^3)$~\cite{BrunerieJames}.
\item The Cayley--Dickson construction and the H-space structure on $S^3$~\cite{BuchholtzCayleyDickson}.
\item The join construction~\cite{RijkeJoin}.
\item Cellular cohomology~\cite{BuchholtzCellular}.
\item Real projective spaces $\mathbb{R}P^n$~\cite{BuchholtzRPn}.
\item Homology theory~\cite{GrahamSyntheticHomology}.
\item Blakers--Massey connectivity theorem~\cite{FavoniaBlakersMassey}.
\item Seifert--van Kampen theorem~\cite{FavoniaSeifertVanKampen}.
\item Nilpotent types and fracture squares~\cite{ScoccolaNilpotent}.
\item Eilenberg--MacLane spaces $K(G,n)$~\cite{LicataEilenbergMacLane}.
\item Spectral sequences~\cite{vanDoornSpectral}.
\item Long exact sequence of homotopy $n$-groups~\cite{BuchholtzLongExact}.
\item Brouwer’s fixed-point theorem in real-cohesion~\cite{ShulmanBrouwer}.
\item Cartan Geometry using a modality~\cite{CherubiniCartan}.
\item Synthetic Cohomology Theory in Cubical Agda~\cite{BrunerieSyntheticCohomology}.
\item Computing Cohomology Rings in Cubical Agda~\cite{LamiauxCohomologyRings}.
\end{itemize}

\subsection{Open Problems}

Some notable open problems and challenges in the formalization of homotopy theory include:

\begin{itemize}
\item Do more “ordinary” ($\leq 2$-truncated) mathematics: Is univalence a practical advantage for formalizing scheme or manifold theory?
\item Prove Freyd’s Adjoint Functor Theorem constructively.
\item Calculate more homotopy groups of spheres.
\item Is $\Sigma A$ a 1-type for all sets $A$?
\item Prove that all spheres $S^n$ are $\infty$-truncated.
\item Define the type $\mathrm{BSU}(2)$ (a delooping of $S^3$); then define $\mathrm{BSU}(n)$, $\mathrm{BSO}(n)$, etc.
\item Define the type of $U$-small semi-simplicial types.
\item For which $0$- or $1$-categories $D$ can we define the type of diagrams $\mathrm{Fun}(D,U)$?
\item Define the type of $(\infty,1)$-categories and develop $(\infty,1)$-category theory.
\item Define the localization type of a $1$-category $C$: the universal type with a functor from $C$.
\item Give a general theory of $\infty$-quotients.
\item Formalize the meta-theory of type theory in itself.
\end{itemize}

\subsection{Discussion}

\begin{itemize}
\item What is a good foundational framework for mathematics? Which criteria are essential, and which are “nice to have”?
\item How does set theory fare according to these criteria?
\item How does (homotopy) type theory fare?
\end{itemize}

Anders aims to provide an efficient and educational environment for exploring and extending these formalizations, particularly through its modal and infinitesimal primitives that support synthetic differential geometry and higher categorical constructions.

\newpage
\begin{thebibliography}{20}

\bibitem{deRham}
Felix Cherubini.
\newblock Cartan Geometry in Modal Homotopy Type Theory.
\newblock \url{https://arxiv.org/pdf/1806.05966.pdf}, 2019.

\bibitem{CHM}
Thierry Coquand, Simon Huber, and Anders Mörtberg.
\newblock On Higher Inductive Types in Cubical Type Theory.
\newblock \url{https://staff.math.su.se/anders.mortberg/papers/cubicalhits.pdf}, 2017.

\bibitem{Huber}
Simon Huber.
\newblock On Higher Inductive Types in Cubical Type Theory.
\newblock \url{http://www.cse.chalmers.se/~simonhu/misc/hcomp.pdf}, 2017.

\bibitem{MLTT72}
Martin-Löf.
\newblock An Intuitionistic Theory of Types, 1972.

\bibitem{MLTT75}
Martin-Löf.
\newblock An Intuitionistic Theory of Types: Predicative Part, 1975.

\bibitem{MLTT80}
Martin-Löf.
\newblock Intuitionistic Type Theory.
\newblock \url{https://raw.githubusercontent.com/michaelt/martin-lof/master/pdfs/Bibliopolis-Book-retypeset-1984.pdf}, 1980.

\bibitem{CIC}
Christine Paulin-Mohring.
\newblock Introduction to the Calculus of Inductive Constructions.
\newblock \url{https://hal.inria.fr/hal-01094195/document}, 2015.

\bibitem{HTS}
Vladimir Voevodsky.
\newblock A simple type system with two identity types.
\newblock \url{https://www.math.ias.edu/vladimir/sites/math.ias.edu.vladimir/files/HTS.pdf}, 2013.

\bibitem{2LTT}
Danil Annenkov, Paolo Capriotti, Nicolai Kraus, and Christian Sattler.
\newblock Two-Level Type Theory and Applications.
\newblock \url{https://arxiv.org/pdf/1705.03307.pdf}, 2019.

\bibitem{CubicalAgda}
Andrea Vezzosi, Anders Mörtberg, and Andreas Abel.
\newblock Cubical Agda: A Dependently Typed Programming Language with Univalence and Higher Inductive Types.
\newblock \url{https://staff.math.su.se/anders.mortberg/papers/cubicalagda.pdf}, 2019.

\bibitem{CCHM}
Cyril Cohen, Thierry Coquand, Simon Huber, and Anders Mörtberg.
\newblock Cubical Type Theory: A Constructive Interpretation of the Univalence Axiom.
\newblock \url{https://staff.math.su.se/anders.mortberg/papers/cubicalagda.pdf}, 2018.

\bibitem{Awodey12}
Steve Awodey.
\newblock Type Theory and Homotopy.
\newblock In \emph{Epistemology versus Ontology: Essays on the Philosophy and Foundations of Mathematics in Honour of Per Martin-Löf}, pages 183--201. Springer, 2012.

\bibitem{Coquand18}
Thierry Coquand.
\newblock A Survey of Constructive Models of Univalence.
\newblock Preprint or Lecture Notes, 2018.

\bibitem{Huber17}
Simon Huber.
\newblock Canonicity for Cubical Type Theory.
\newblock \emph{Journal of Automated Reasoning}, 61(1-4):173--205, 2017.

\bibitem{Streicher91}
Thomas Streicher.
\newblock \emph{Semantics of Type Theory: Correctness, Completeness and Independence Results}.
\newblock Birkhäuser, Basel, 1991.

\bibitem{Bezem14}
Marc Bezem, Thierry Coquand, and Simon Huber.
\newblock A Model of Type Theory in Cubical Sets.
\newblock Preprint, arXiv:1406.1731, 2014.

\bibitem{Shulman15}
Michael Shulman.
\newblock Univalence for Inverse Diagrams and Homotopy Canonicity.
\newblock \emph{Mathematical Structures in Computer Science}, 25(5):1203--1277, 2015.

\bibitem{Hofmann97}
Martin Hofmann.
\newblock Syntax and Semantics of Dependent Types.
\newblock In \emph{Semantics and Logics of Computation}, pages 79--130. Cambridge University Press, 1997.

\bibitem{ThreeHIT}
Andrej Bauer and Niels van der Weide.
\newblock The Three-HITs Theorem.
\newblock DOI:10.4230/LIPIcs.CVIT.2016.23, \url{https://5ht.co/three.pdf}.

% ==================== HoTT Mathematics ====================

\bibitem{SymmetryBook}
Marc Bezem, Ulrik Buchholtz, Bjorn Dundas, Dan Grayson, et al.
\newblock Symmetry: A book project about group theory from the univalent point of view.
\newblock 2020. \url{https://github.com/UniMath/SymmetryBook}.

\bibitem{BrunerieJames}
Guillaume Brunerie.
\newblock The James Construction and $\pi_4(S^3)$ in Homotopy Type Theory.
\newblock \emph{Journal of Automated Reasoning}, 63.2, pp. 255--284, 2019.
\newblock DOI: \url{https://doi.org/10.1007/s10817-018-9468-2}, arXiv:1710.10307.

\bibitem{BuchholtzHigherGroups}
Ulrik Buchholtz, Floris van Doorn, and Egbert Rijke.
\newblock Higher Groups in Homotopy Type Theory.
\newblock In \emph{Proceedings of the 33rd Annual ACM/IEEE Symposium on Logic in Computer Science (LICS '18)}, pp. 205--214, 2018.
\newblock DOI: \url{https://doi.org/10.1145/3209108.3209150}, arXiv:1802.04315.

\bibitem{BuchholtzCellular}
Ulrik Buchholtz and Favonia.
\newblock Cellular Cohomology in Homotopy Type Theory.
\newblock In \emph{Proceedings of the 33rd Annual ACM/IEEE Symposium on Logic in Computer Science (LICS '18)}, pp. 521--529, 2018.
\newblock DOI: \url{https://doi.org/10.1145/3209108.3209188}, arXiv:1802.02191.

\bibitem{BuchholtzCayleyDickson}
Ulrik Buchholtz and Egbert Rijke.
\newblock The Cayley--Dickson construction in homotopy type theory.
\newblock \emph{Higher Structures} 2.1, pp. 30--41, 2018.
\newblock arXiv:1610.01134.

\bibitem{BuchholtzLongExact}
Ulrik Buchholtz and Egbert Rijke.
\newblock The long exact sequence of homotopy $n$-groups.
\newblock 2019. arXiv:1912.08696.

\bibitem{BuchholtzRPn}
Ulrik Buchholtz and Egbert Rijke.
\newblock The real projective spaces in homotopy type theory.
\newblock In \emph{32nd Annual ACM/IEEE Symposium on Logic in Computer Science (LICS)}, pp. 1--8, 2017.
\newblock DOI: \url{https://doi.org/10.1109/LICS.2017.8005146}, arXiv:1704.05770.

\bibitem{ChristensenLocalization}
J. Daniel Christensen, Morgan Opie, Egbert Rijke, and Luis Scoccola.
\newblock Localization in Homotopy Type Theory.
\newblock 2018. arXiv:1807.04155.

\bibitem{vanDoornSpectral}
Floris van Doorn, Jeremy Avigad, Steve Awodey, Ulrik Buchholtz, Egbert Rijke, and Michael Shulman.
\newblock Spectral Sequences in Homotopy Type Theory.
\newblock 2019. \url{https://github.com/cmu-phil/Spectral}.

\bibitem{FavoniaBlakersMassey}
Favonia, Eric Finster, Daniel R. Licata, and Peter L. Lumsdaine.
\newblock A Mechanization of the Blakers--Massey Connectivity Theorem in Homotopy Type Theory.
\newblock In \emph{31st Annual ACM/IEEE Symposium on Logic in Computer Science (LICS)}, pp. 1--10, 2016.
\newblock arXiv:1605.03227.

\bibitem{FavoniaSeifertVanKampen}
Favonia and Michael Shulman.
\newblock The Seifert--van Kampen Theorem in Homotopy Type Theory.
\newblock In \emph{25th EACSL Annual Conference on Computer Science Logic (CSL 2016)}, LIPIcs vol. 62, pp. 22:1--22:16, 2016.
\newblock DOI: \url{https://doi.org/10.4230/LIPIcs.CSL.2016.22}.

\bibitem{GrahamSyntheticHomology}
Robert Graham.
\newblock Synthetic Homology in Homotopy Type Theory.
\newblock 2017. arXiv:1706.01540.

\bibitem{KrausFreeHigherGroups}
Nicolai Kraus and Thorsten Altenkirch.
\newblock Free Higher Groups in Homotopy Type Theory.
\newblock In \emph{Proceedings of the 33rd Annual ACM/IEEE Symposium on Logic in Computer Science (LICS '18)}, pp. 599--608, 2018.
\newblock DOI: \url{https://doi.org/10.1145/3209108.3209183}, arXiv:1805.02069.

\bibitem{LicataEilenbergMacLane}
Daniel R. Licata and Eric Finster.
\newblock Eilenberg--MacLane Spaces in Homotopy Type Theory.
\newblock In \emph{CSL-LICS '14}, 2014.
\newblock DOI: \url{https://doi.org/10.1145/2603088.2603153}.

\bibitem{RijkeJoin}
Egbert Rijke.
\newblock The join construction.
\newblock 2017. arXiv:1701.07538.

\bibitem{RijkeModalities}
Egbert Rijke, Michael Shulman, and Bas Spitters.
\newblock Modalities in homotopy type theory.
\newblock \emph{Logical Methods in Computer Science} 16.1, 2020.
\newblock arXiv:1706.07526, \url{https://lmcs.episciences.org/6015}.

\bibitem{ScoccolaNilpotent}
Luis Scoccola.
\newblock Nilpotent Types and Fracture Squares in Homotopy Type Theory.
\newblock 2019. arXiv:1903.03245.

\bibitem{ShulmanBrouwer}
Michael Shulman.
\newblock Brouwer’s fixed-point theorem in real-cohesive homotopy type theory.
\newblock \emph{Mathematical Structures in Computer Science} 28.6, pp. 856--941, 2018.
\newblock DOI: \url{https://doi.org/10.1017/S0960129517000147}, arXiv:1509.07584.

\bibitem{CherubiniCartan}
Felix Wellen (Cherubini).
\newblock Cartan Geometry in Modal Homotopy Type Theory.
\newblock 2018. arXiv:1806.05966.

\bibitem{BrunerieSyntheticCohomology}
Guillaume Brunerie, Axel Ljungström, and Anders Mörtberg.
\newblock Synthetic Cohomology Theory in Cubical Agda.
\newblock \url{https://staff.math.su.se/anders.mortberg/papers/zcohomology.pdf}.

\bibitem{LamiauxCohomologyRings}
Thomas Lamiaux, Axel Ljungström, and Anders Mörtberg.
\newblock Computing Cohomology Rings in Cubical Agda.
\newblock arXiv:2212.04182, \url{https://hott-uf.github.io/2023/slides/Lamiaux.pdf}.

\end{thebibliography}

\end{document}
